%----------------------------------------------------------------------------------------
%				Creato da Mich
%   Updated by Simone Pessotto 04/08/2015
%   Updated by Riccardo Montagnin 13/04/2017
%   Updated by Federico Silvio Busetto 15/09/2017
%----------------------------------------------------------------------------------------

\documentclass[11pt,a4paper]{article}

\usepackage[utf8x]{inputenc}
\usepackage{helvet}
\usepackage{multirow}
\renewcommand*\familydefault{\sfdefault}
\usepackage[T1]{fontenc}
\linespread{1.2}
\usepackage{fancyhdr}
\usepackage{hyperref}
\usepackage{lastpage}
\usepackage{amsfonts, amsmath, amsthm, amssymb}
\usepackage{graphicx}

\usepackage{tabularx}
\newcommand{\specialcell}[2][c]{%
    \begin{tabular}[#1]{@{}l@{}}#2
    \end{tabular}
}

% Specifies custom hyphenation points for words or words that shouldn't be hyphenated at all
\hyphenation{ionto-pho-re-tic iso-tro-pic fortran} 
\hypersetup{
    colorlinks=true,
    linkcolor=black,
    urlcolor=blue
}

\usepackage[
a4paper,
top=2.5cm,
bottom=2.5cm,
left=1.5cm,
right=1.5cm,
head=30pt
]{geometry}

\usepackage[italian]{isodate}
\usepackage{arydshln}

%----------------------------------------------------------------------------------------
%	PAGE STYLE
% ---------------------------------------------------------------------------------------
\pagestyle{fancy}
\fancyhf{}
\setlength{\headheight}{2cm} 

% Delete the paragraph indentation
\setlength{\parindent}{0pt}

%----------------------------------------------------------------------------------------
% 	HEADING STILE
%----------------------------------------------------------------------------------------
\newcommand\textline[4][t]{%
    \noindent\parbox[#1]{.333\textwidth}{\raisebox{-0.40\height}{#2}}%
    \parbox[#1]{.333\textwidth}{\centering #3}%
    \parbox[#1]{.333\textwidth}{\raggedleft #4}%
}

\lhead{
    \textline[t]{\includegraphics[width=1cm, keepaspectratio=true]{img/UniPd.png}}
    {Piano di lavoro stage presso \\\ragioneSocAzienda}
    {\nomeStudente \cognomeStudente\\ (\matricolaStudente)}
}

% Line under the heading
\renewcommand{\headrulewidth}{0.4pt}  

%-----------------------------------------------------------------------------------------
%   FOOTER STYLE
%-----------------------------------------------------------------------------------------

%***PIÈ DI PAGINA***
%\lfoot{\includegraphics[keepaspectratio = true, width = 25px] {img/UniPd.png} \textit{\nomeStudente 
%\cognomeStudente (\matricolaStudente)}\\ \footnotesize{\emailStudente}}
\rfoot{\thepage/\pageref{LastPage}} %per le prime pagine: mostra solo il numero romano
\cfoot{}
\renewcommand{\footrulewidth}{0.4pt}   %Linea sopra il piè di pagina